\subsection{The radial-loop class of every global cap}

\begin{proposition}[B.4.32; Proved; vanishing of the global affine loop defect]
\label{push:B432}
Suppose a proper complex-analytic filling of a finite parameter
cover of the actual six-sphere family exists. Make the
modifications in B.4.31 and delete its finite exceptional set
from the marked base. Use the ordinary local origins of
B.4.19 as reference. There is a well-defined normalized
radial-loop defect of a local extending cap section,
\[
 \lambda\in
 V:=\bigl(\Lambda\otimes\Q\bigr)/\Q\widehat\delta .
 \tag{B.4.32.1}
\]
It is independent of the section, of the finite parameter root
needed to obtain that section, and of the modifications used.
It patches to a flat section of the actual local system $V$
over the remaining ordinary base, and this section is zero.

In particular, let $t$ be the trace of an extending finite
multisection of degree $d$, and let $o$ be the reference ordinary
origin. The radial-loop defect of $t-do$ in $V$ is zero.
The same holds for an origin obtained from multiplied such
traces by bounded logarithmic differences along $r=0$.
This conclusion controls a loop class; it does not assert
that $t-do$ is bounded or extends holomorphically at $r=0$.
\end{proposition}

\begin{proof}
After a common further parameter cover we may compare with
$w=r^6$ in B.4.19; further positive powers are allowed throughout.
Retain and resolve the reduced main component of the putative
filling. The extension of the marked map uses B.4.4, whose
smooth fourfold proof requires no condition on the central
divisor. B.4.31 supplies a proper smooth marked graph with
normal-crossings central support, possibly with multiplicities,
and the integral meridian conclusion. These operations are
the identity on the punctured family.

\emph{The connecting class and its normalization.}
On a simply connected ordinary base chart, write
\[
 T_r=I-mE,\qquad
 E=\widehat\delta\otimes\gamma,\qquad m>0,\qquad E^2=0.
 \tag{B.4.32.2}
\]
For a holomorphic section $a$ of the Jacobian on the punctured
parameter disc, its connecting class for the Lie-bundle
exponential sequence belongs to
\[
 H^1(\Delta_r^*,\Lambda)
       =\Lambda/(T_r-I)\Lambda .
\]
Equivalently, continue a local additive lift around the positive
loop and record its endpoint period. A change of lift alters
the period label by $(T_r-I)\ell$. Since this lies in
$\Z\widehat\delta$, the image
\[
 \nu_r(a)\in V
 \tag{B.4.32.3}
\]
is well defined. No single-valued additive lift or meromorphic
extension of $a$ has been assumed. In a holomorphic family of
sections this class is locally constant in the ordinary base:
the integral period label is continuous and discrete.

Under $r=R^e$, restriction of the connecting class sends a
period label $v$ to
\[
 \sum_{j=0}^{e-1}T_r^jv
   =ev-\frac{me(e-1)}2Ev
   \equiv ev\pmod{\Q\widehat\delta}.
 \tag{B.4.32.4}
\]
If $b$ is a cap section on this root cover, define
\[
 \lambda(b;o)=\frac1e\,\nu_R\bigl(b-o(R^e)\bigr).
 \tag{B.4.32.5}
\]
Formula (B.4.32.4) proves invariance under further roots and
fixes the normalization used below for finite traces.

\emph{The actual affine changes of reference origin.}
This step uses the specified origins, not arbitrary zeros of
abstract isomorphic tori. Retain
\[
 v=\frac{\log r}{2\pi i},\qquad c=-12v,\qquad
 k_i=\frac{\log K}{2\pi i},\qquad
 K(x)=\frac{x^4}{(x-1)^3}.
\]
The ordinary gauge of B.4.19 is
\[
 \widetilde\zeta_i=\zeta+k_i v e,\qquad
 e=(0,1)^t=\Pi_c\widehat\delta .
 \tag{B.4.32.6}
\]
Its zero is the local reference origin. On an ordinary overlap
$k_j-k_i=n_{ij}\in\Z$. The extra translation is $n_{ij}ve$;
its radial-loop label is $n_{ij}\widehat\delta$, hence zero
in $V$.

The linear Lie-frame changes are independent of $c$ and $r$.
For the two generating finite peripheral changes they are
\[
 R_1=
 \begin{pmatrix}-1/\tau&0\\(1-\mu)/\tau&1\end{pmatrix},
 \qquad
 R_2=
 \begin{pmatrix}1/\tau&0\\-\mu/\tau&1\end{pmatrix}.
 \tag{B.4.32.7}
\]
They satisfy $R_je=e$ and $R_j\Pi_c=\Pi'_cA_j$, with
source and target arguments retained as in B.1.2. Their
compositions therefore act on (B.4.32.3) by the actual lattice
transport and have no additional parameter-dependent
Lie-frame contribution.

The finite-root connector is the potentially nonconstant
affine term. Its full cancellation is B.4.19.4:
\[
 \begin{split}
  \sigma_{j,c}&=w_j\Pi_c v_j,\qquad
   k-k_j=12\epsilon_jw_j,\qquad
   \gamma(v_j)=\epsilon_j,\\
  \sigma_{j,c}+(k-k_j)ve
     &=w_j\Pi_{\rm part}v_j .
 \end{split}
 \tag{B.4.32.8}
\]
The last expression is independent of $r$, since $c=-12v$.
It has zero radial-loop defect. The remaining finite deck
maps have the multiplicative factors $r^4$ and $r^{-3}$
in B.1.2, which add only integral $\widehat\delta$ labels.
Changing a sector's period representative adds
$(T_r-I)\ell$, again in $\Q\widehat\delta$.

On a punctured cusp sector the comparison is, by B.4.19.5,
\[
 G(s)+6ve,\qquad
 G(s)=\left(\frac{\mu s}{2},
                 \frac{Cs}{12}-\frac{s^2}{24}\right)^t .
 \tag{B.4.32.9}
\]
Here $G$ is independent of $r$. Its additional radial-loop
label is $6\widehat\delta$. A shift by all twelve components
need not be the identity on the ruled normalization; it may
translate the elliptic base by $6P$. As an actual period
change its radial-loop contribution is nevertheless
$(T_r-I)\widehat\gamma$, which is in the same $\delta$ line.

These equations prove that the reference origins determine one
globally consistent loop class in $V$. They do not assert that
the origins themselves glue. A different coherent closing
related by a global scalar $\C^*$ gauge has only an integral
$\widehat\delta$ winding and gives the same class. A further
power of the parameter multiplies the displayed loop labels
and leaves all these conclusions unchanged.

\emph{Comparison of sections extending into one cap.}
Choose an ordinary base disc avoiding the finite set in
B.4.31. Its closed branch strata map submersively to the
disc. The local finite root and compatible toric subdivision
argument of B.4.30 gives a proper semistable model with this
ordinary base parameter retained. Only a model over this
small disc is required. On an overlap of two choices use a
common further root and common refinement. No global
semistable-reduction theorem over the entire punctured
$\PP^1$ is used.

The central surface fibres of this model are connected.
Indeed the proper map from its normal dominating total space
to the normal surface base has a finite Stein factor.
It is generically degree one because the ordinary torus
fibres are connected. A finite birational map to a normal
base is an isomorphism.

The saturation argument (B.4.30.4)--(B.4.30.5) still applies.
It does not require rank-two monodromy: it only requires the
known normalized rank-two form subbundle of the canonical
log-de Rham extension. The ordinary period matrix supplies
that subbundle, since the $r$ logarithm is solely in the
$\widehat\delta\otimes\gamma$ term and the $w,\delta$
columns stay the identity. The analytic comparison theorem
used here is precisely
\citet[Theorem 6.3]{IKN2005E}, with the hypotheses and erratum
scope already checked in B.4.30.

Rank two in B.4.30 supplied the integral local residues.
Here B.4.31 supplies what is needed more directly: every
local angular class is an integral multiple of
$\widehat\delta$. Thus the first normalized form has residue
zero and the second has integral residue. The convergent
primitive calculation (B.4.30.6) gives exponentials in
\[
              \OO^*\times M_r^{\rm gp}.
 \tag{B.4.32.10}
\]
This is the exact bounded character sheaf along the strict
$r$ axis. In particular the first character is a holomorphic
unit across $r=0$ and has zero radial winding. The second
is a power of $r$ times a unit and contributes only an
integral $\widehat\delta$ winding.

Given two extending sections, join their limits by the finite
chart chain of B.4.30. The central connectedness just proved
allows different components as endpoints. The resulting
character comparison lies in (B.4.32.10), so
\[
               \nu_R(b-b')=0\quad\hbox{in }V.
 \tag{B.4.32.11}
\]
Different roots are compared on a common further root.
Proper lifts through refinements are meromorphic on the
two-dimensional section base and holomorphic away from
isolated points. Comparing on the nonempty open set where
the lifts are holomorphic suffices: both loop classes are
locally constant and hence the equality holds throughout
the ordinary overlap. Thus the normalized class is
independent of all the choices.

Smooth points of a reduced central fibre have transverse
holomorphic sections, so this common class exists locally.
If an extending finite multisection of degree $d$ is used
instead, split it after a finite cyclic root at a generic
ordinary value. Its $d$ branches have the same normalized
class. Additivity of the connecting homomorphism and
(B.4.32.4) give
\[
            \nu_r(t-do)=d\,\lambda .
 \tag{B.4.32.12}
\]
This also proves independence from finite trace
presentations without presuming an extending chosen origin.

\emph{Global transport and its fixed space.}
On one ordinary base chart choose one extending section.
Its normalized integral loop label is locally constant.
On overlaps, (B.4.32.11) removes the choice of cap section,
and (B.4.32.6)--(B.4.32.9) remove the reference-origin
ambiguity in $V$. The local classes therefore patch under
the actual homology transport. Hence $\lambda$ is a flat
section of $V$ over the ordinary base with finitely many
additional points removed.

In the quotient basis
$(\widehat\gamma,\widehat u,\widehat w)$, the matrices are
\[
 \overline A_1=
 \begin{pmatrix}1&0&0\\6&0&1\\-6&-1&-1\end{pmatrix},
 \qquad
 \overline A_2=
 \begin{pmatrix}1&0&0\\0&0&-1\\-6&1&0\end{pmatrix}.
 \tag{B.4.32.13}
\]
Their rational fixed spaces are respectively
\[
 \Q(1,2,-4),\qquad \Q(1,3,-3),
 \tag{B.4.32.14}
\]
whose intersection is zero. Removing additional ordinary
base values does not decrease the monodromy image: the
map on fundamental groups to the three-punctured sphere
is surjective. Therefore $\lambda=0$.

Even if one retained only a finite orbit of trace branches,
the conclusion would be the same. The covector $\gamma$
is invariant. Averaging a finite orbit would give a common
invariant vector with its original $\gamma$ coefficient,
so that coefficient is zero. The remaining orbit lies in
the elliptic two-dimensional subspace, where the two
unipotents in B.4.31 have distinct fixed lines and no
nonzero finite orbit.

Finally, a bounded logarithmic difference on the strict
$r$ axis has the character form (B.4.32.10), hence its
loop label lies in $\Q\widehat\delta$. Adding such a
difference to a multiplied finite trace preserves its
zero defect. This proves the last assertion and in
particular gives zero after applying the invariant
covector $\gamma$.
\end{proof}

\begin{remark}[Scope of the loop calculation]
No assertion of a bounded, meromorphic or holomorphic
comparison with the reference origin has been made.
Essential translations with zero loop defect remain
possible. The conclusion uses the global marked
$\PP^1$ base and its two peripheral matrices; an isolated
rank-one degeneration over an ordinary base disc has
no such constraint. Arbitrary singular proper fillings
are covered by resolving their reduced main component
before B.4.4 and B.4.31; their central multiplicities
are retained until the local finite-root calculations.
\end{remark}

The exact quotient fixed spaces, the root-transfer identity,
the finite connector cancellation and the cusp loop label
are checked in
\path{code/d_push_global_loop_defect.py}, with its saved
output and JSON record. The proof of the global geometric
comparison is the argument above.
